\chapter{Nghe Giảng Như Nghe Gió}
\chapterintro{Âm thanh đi qua rất đủ, chỉ tiếc ý nghĩa không ở lại nhiều}{Đây là cấp độ nâng cao của việc không tập trung: nghe mà như không.}

\begin{lucbatbox}{Lục bát: Gió qua tai nhẹ}
Thầy cô giảng suốt một hồi\
Tai em nghe đủ, đầu thời không lưu\
Chữ đi như gió cuối chiều\
Chạm qua rất nhẹ rồi vèo mất tăm\
Tới khi mở vở hỏi thăm\
Mới hay bài giảng đi ngang cuộc đời
\end{lucbatbox}

\begin{tutuyetbox}{Thất ngôn tứ tuyệt: Âm thanh có, ý không}
Lời giảng vẫn tới tai em\
Nhưng như đi lạc qua miền mây xa\
Nghe xong tưởng đã hiểu mà\
Đụng vào chi tiết mới ra mịt mù
\end{tutuyetbox}

\begin{ghichu}
Nghe giảng như nghe gió là tình trạng nguy hiểm ở chỗ người trong cuộc vẫn tưởng mình đã nghe đủ rồi.
\end{ghichu}

\ngancach

\begin{lucbatbox}{Lục bát: Gật gù cho có}
Đoạn nào thầy nhấn hơi sâu\
Em còn gật gật như cầu rất thông\
Bạn nhìn tưởng đã hiểu xong\
Chỉ mình em biết trong lòng đang trôi\
Nếu mà dừng lại một hồi\
Hỏi thêm một chút chắc rồi khá hơn
\end{lucbatbox}

\begin{tutuyetbox}{Thất ngôn tứ tuyệt: Trôi qua rất êm}
Tiết học trôi rất nhẹ nhàng\
Nhẹ đến mức cả bài sang mất rồi\
Em còn ngồi đó mà thôi\
Kiến thức thì đã ra ngoài từ lâu
\end{tutuyetbox}